\documentclass[10pt,a4paper]{amsart}
\usepackage[margin=19mm]{geometry}
\usepackage{amsmath,amssymb,mathtools}
\usepackage[scr=rsfs,cal=euler]{mathalfa}
\usepackage{xfrac}
\usepackage[unicode,hidelinks,bookmarks=false]{hyperref}
\newcommand{\ie}{i.e.\ }
\newcommand{\cf}{cf.\ }
\newcommand{\resp}{respectively\ }
\newcommand{\Subsection}{\subsection}
\providecommand{\nobreakdash}{\nobreak\hbox{-}}
\setlength{\emergencystretch}{2em}
\multlinegap0pt
\usepackage{tikz}
\usepackage{amssymb}
\usepackage{mathtools}
\usepackage{float}
\newtheorem{proposition}{Proposition}
\newtheorem{lemma}{Lemma}
\def\Li{\operatorname{Li}}
\def\Log{\operatorname{Log}}
\renewcommand{\Re}{{\rm Re}}
\title{Inequalities for the number of $t$-hooks in two partition classes arising from sum-product identities}
\author{Aritram Dhar, Byungchan Kim, Eunmi Kim, Ae Ja Yee}
\date{Source: arXiv:2603.01112v1 (CC BY 4.0)}
\begin{document}
\maketitle

\noindent\emph{Source and licence.} Inequalities for the number of $t$-hooks in two partition classes arising from sum-product identities. Version arXiv:2603.01112v1 (CC BY 4.0), distributed by arXiv under the Creative Commons Attribution 4.0 International licence, http://creativecommons.org/licenses/by/4.0/, as recorded in the arXiv licence metadata for this exact version and shown on the abstract page https://arxiv.org/abs/2603.01112v1. Full licence notice: \texttt{licenses/arxiv-2603.01112-CC-BY-4.0.txt}.\par\smallskip

\noindent\textit{Editorial scope.} Counted unit: the article's proposition prop:S11\_asymp (source0.tex lines 914--919). The article's complete original proof of it is delivered with it (source0.tex lines 920--985, one \texttt{proof} environment of 66 lines). No mathematics is written, repaired or re-derived: the statement and the proof are the article's own lines, in the article's own notation, and no retained line is substituted or re-flowed.  Retained uncounted as the source-local context the proof needs: lines 619--624; lines 704--709; lines 782--796; lines 849--858.  Nothing was omitted from any retained span, so the package carries no omission record.  No retained line contains a citation, so the wrapper carries no bibliography and no bibliography entry.  No retained line contains a figure environment, an image-inclusion command or drawing code, so no image is delivered and the package declares no asset dependency.  Editorial formatting only, in addition: an AMS \texttt{amsart} preamble that restates the article's own declarations and macros used by the retained lines, namely the environments \texttt{proposition}, \texttt{lemma} on separate counters, together with the standard packages those lines need.  The retained environments print in this document as Proposition 1, Lemma 1 in order of appearance, whereas the article prints the same items with the numbering of its own full document.  The complete native source of the article as distributed in the arXiv e-print for this exact version is bundled unchanged as \texttt{sources/195540/source0.tex}.
\begin{proposition}\label{prop:S11_narrow}
	If $y \ll \varepsilon^{\frac{2}{5}+\delta}$ for some $\delta>0$, then we have, as $z \to 0$, 
	\[
		\mathcal{S}_{1,1}(q) = \sum_{n \in \mathcal{N}_{\varepsilon}} \frac{nq^{n^2}}{(q;q)_n} = \frac{\sqrt{\phi} }{5^{1/4}}\log(\phi)  z^{-1} e^{\frac{\pi^2}{15 z}} + O\left(\varepsilon^{-\frac{9}{10}} e^{\frac{\pi^2}{15\varepsilon}}\right).
	\]
\end{proposition}

\begin{proposition}\label{prop:far}
	For every $L \in \mathbb{N}$, as $\Re(z) \to 0$ in the right half-plane, then
	\[
		\left| \mathcal{E}_{1,1}(q) \right| = \left|\sum_{n \in \mathbb{N} \setminus \mathcal{N}_{\varepsilon}}  \frac{nq^{n^2}}{(q;q)_n} \right| \ll \varepsilon^{L-1} e^{\frac{\pi^2}{15\varepsilon}}.
	\]
\end{proposition}

\begin{proposition}\label{prop:wider}
	Let $n \in \mathcal{N}_{\varepsilon}$ and $v \in \mathbb{C}$ be such that
	\[
		n = \frac{\log(\phi)}{\varepsilon} + \frac{v}{\sqrt{z}}. 
	\]
	If $y \ll 1$, then we have, uniformly in $v$,
	\begin{multline*}
		\frac{n q^{n^2}}{(q;q)_n} = \frac{\log(w^{-1})}{\sqrt{2\pi z(1-w)}} \exp\left(\frac{\Lambda (y)}{z} + \Log\left(\frac{w^2}{1-w}\right)\frac{v}{\sqrt{z}} - \frac{2-w}{2(1-w)}  v^2   \right) \\
		+ O\left(\varepsilon^{\frac{1}{5}}\right) \frac{1}{\sqrt{z}} \exp\left(\frac{\Lambda (y)}{z} + \Log\left(\frac{w^2}{1-w}\right)\frac{v}{\sqrt{z}} - \frac{2-w}{2(1-w)}  v^2   \right),
	\end{multline*}
	where $w := \phi^{-1-iy}$ and
	\begin{equation}\label{eq:Lamma}
		\Lambda(y) := \frac{\pi^2}{6} - \log^2(\phi) (1+iy)^2 - \Li_2 ( \phi^{-(1+iy)} ).
	\end{equation}
\end{proposition}

\begin{lemma}\label{lem:s}
	Let $s(y):=\Re\left(\frac{\Lambda(y)}{1+iy} - \frac{\pi^2}{15}\right)$. Then we have the following.
	\begin{enumerate}
		\item $s(y) \leq 0$ for all $y \in \mathbb{R}$, and the equality holds if and only if $y=0$.
		\item As $y \to 0$,
		\[
			s(y) = - \left(\frac{\pi^2}{15} + \left( 3- \frac{\phi}{2}\right)\log^2(\phi) \right)y^2 + O\left(y^4\right).
		\]
	\end{enumerate}
\end{lemma}

\begin{proposition}\label{prop:S11_asymp}
	If $y \ll 1$, then we have that as $z \to 0$,
	\[
		S_{1,1}(q) = \frac{\phi^{\frac12} }{5^{\frac14}}\log(\phi)  z^{-1} e^{\frac{\pi^2}{15 z}} +  O\left(\varepsilon^{-\frac{9}{10}} e^{\frac{\pi^2}{15\varepsilon}}\right).
	\]
\end{proposition}

\begin{proof}
By Proposition \ref{prop:far}, we have, for all $L \in \mathbb{N}$, as $z\to 0$,
\[
	\left|\sum_{n \in \mathbb{N} \setminus \mathcal{N}_{\varepsilon}}  \frac{nq^{n^2}}{(q;q)_n} \right| \ll \varepsilon^{L-1} e^{\frac{\pi^2}{15\varepsilon}}.
\]
Thus, the desired asymptotic follows from the following claim:
\[
	z e^{-\frac{\pi^2}{15\varepsilon}} \left(\sum_{n \in \mathcal{N}_{\varepsilon}} \frac{nq^{n^2}}{(q;q)_n} - \frac{\phi^{\frac12} }{5^{\frac14}}\log(\phi)  z^{-1} e^{\frac{\pi^2}{15 z}} \right) \ll \varepsilon^{\frac{1}{10}} .
\]
We verify this claim by splitting the range of $y$ into two pieces, $y \ll \varepsilon^{\frac25+\delta}$ and  $\varepsilon^{\frac12-\delta}\ll y \ll 1$. By choosing $\delta >0$ sufficiently small, these two pieces can cover all necessary range $y \ll 1$. We first note that for $y \ll \varepsilon^{\frac25+\delta}$ for some $\delta>0$, then, by Proposition \ref{prop:S11_narrow}, we have
\[
	\sum_{n \in \mathcal{N}_{\varepsilon}} \frac{nq^{n^2}}{(q;q)_n} - \frac{\phi^{\frac12} }{5^{\frac14}}\log(\phi)  z^{-1} e^{\frac{\pi^2}{15 z}} \ll \varepsilon^{-\frac{9}{10}} e^{\frac{\pi^2}{15\varepsilon}}.
\]

For the remaining range $\varepsilon^{\frac12-\delta}\ll y \ll 1$ for some $\delta>0$, we will show that, for all $L \in \mathbb{N}$,
\begin{equation}\label{eq:first_part}
	z e^{-\frac{\pi^2}{15\varepsilon}}\sum_{n \in \mathcal{N}_{\varepsilon}} \frac{nq^{n^2}}{(q;q)_n} = O\left(\varepsilon^L\right).
\end{equation}
If $y \ll 1$, by Proposition \ref{prop:wider},
\begin{multline*}
	z e^{-\frac{\pi^2}{15\varepsilon}}\sum_{n \in \mathcal{N}_{\varepsilon}} \frac{nq^{n^2}}{(q;q)_n} = \log(w^{-1})\sqrt{\frac{z}{2\pi(1-w)}} e^{\frac{1}{\varepsilon}\left(\frac{\Lambda(y)}{1+iy}-\frac{\pi^2}{15}\right)} \sum_{n \in \mathcal{N}_{\varepsilon}} e^{-\frac{2-w}{2(1-w)}v^2 +\frac{v}{\sqrt{z}} \Log\left(\frac{w^2}{1-w}\right)}\\
	+ O\left(\varepsilon^{\frac15}\right) \sqrt{z} e^{\frac{1}{\varepsilon}\left(\frac{\Lambda(y)}{1+iy}-\frac{\pi^2}{15}\right)} \sum_{n \in \mathcal{N}_{\varepsilon}}  e^{-\frac{2-w}{2(1-w)}v^2 +\frac{v}{\sqrt{z}} \Log\left(\frac{w^2}{1-w}\right)},
\end{multline*}
where $w=\phi^{-1-iy}$. 
Thus, to prove \eqref{eq:first_part}, it is sufficient to show that the exponent
\begin{equation}\label{eq:exp}
	\frac{1}{\varepsilon}\left(\frac{\Lambda(y)}{1+iy}-\frac{\pi^2}{15}\right) -\frac{2-w}{2(1-w)}v^2 +\frac{v}{\sqrt{z}} \Log\left(\frac{w^2}{1-w}\right)
\end{equation}
has negative real part of size $\gg \varepsilon^{-\delta_0}$ for some $\delta_0>0$ since the other terms are bounded as $z\to 0$. 

If $\varepsilon^{\frac12-\delta} \ll y \ll 1$ for some $\delta>0$, then %$ \delta < \frac12$.
the real part of $\frac{1}{\varepsilon}\left(\frac{\Lambda(y)}{1+iy}-\frac{\pi^2}{15}\right)$ is $\frac{s(y)}{\varepsilon}$, which is negative by Lemma \ref{lem:s} (1). We also obtain, as $\varepsilon \to 0$, 
\[
	\frac{s(y)}{\varepsilon} \gg \frac{y^2}{\varepsilon} \gg \varepsilon^{-2\delta}
\]
by Lemma \ref{lem:s} (2). 
From the Taylor series expansion, as $y \to 0$, 
\[
	\Re\left( \Log\left(\frac{w^2}{1-w}\right) \right) = \log \left|\frac{w^2}{1-w} \right| = -\frac{2+\sqrt{5}}{2} \log^2(\phi) y^2 + O\left(y^4\right) \ll y^2.
\]
As $n \in \mathcal{N}_{\varepsilon}$, we have $\left|\frac{v}{\sqrt{z}}\right| \leq \varepsilon^{-\frac35}$, which implies
\[
	\Re\left( \frac{v}{\sqrt{z}} \Log\left(\frac{w^2}{1-w}\right) \right) \ll \varepsilon^{-\frac35} y^2.
\]	
Since $\frac{y^2}{\varepsilon} \gg \varepsilon^{-\frac35} y^2$, $\frac{s(y)}{\varepsilon}$ dominates $\Re\left( \frac{v}{\sqrt{z}} \Log\left(\frac{w^2}{1-w}\right) \right)$. For the term $\frac{2-w}{2(1-w)}v^2$, we consider the following two cases. 
\begin{enumerate}
	\item Suppose $\varepsilon^{\frac12-\delta} \ll y \ll \varepsilon^{\frac14}$. As $v$ is a real multiple of $\sqrt{z}$ and $y\ll \varepsilon^{\frac14}$, we have that for $y \ll \varepsilon^{\frac14}$,
	\[
		\Re \left(\frac{2-w}{1-w} (1+iy)\right) = \frac{2+\phi^{-2}-3\phi^{-1}\cos\left(y\log(\phi)\right)+y\phi^{-1}\sin\left(y\log(\phi)\right)}{|1-w|^2} >0.
	\]
	Thus, $-\frac{2-w}{2(1-w)}v^2$ has negative real part as $z\to 0$.
	
	\item If $\varepsilon^{\frac13} \ll y \ll 1$, then $|z|\ll \varepsilon$. As $n \in \mathcal{N}_{\varepsilon}$, we have $|v|\ll \varepsilon^{-\frac1{10}}$. Thus, with the fact $\left|\frac{2-w}{2(1-w)}\right| \leq \frac{2-\phi^{-1}}{2(1-\phi^{-1})}$, we find $-\frac{2-w}{2(1-w)}v^2 \ll  \varepsilon^{-\frac15}$. %Since $\delta> \frac14> \frac1{10}$, $\frac{s(y)}{y}$ dominates this term.
\end{enumerate}
In either case, the exponent $\frac{s(y)}{\varepsilon}$ has size $\gg \varepsilon^{-2\delta}$ and dominates the other exponents in \eqref{eq:exp} that have positive real part. Hence, we prove the claim \eqref{eq:first_part} by taking $\delta_0=2\delta$.

Next, in the expression $\frac{\phi^{\frac12} }{5^{\frac14}}\log(\phi) e^{\frac{\pi^2}{15 z}-\frac{\pi^2}{15\varepsilon}}$, the exponential part has negative real part of size
\[
	\Re\left(\frac{\pi^2}{15 z}-\frac{\pi^2}{15\varepsilon}\right)% = \frac{\pi^2}{15\varepsilon} \left( \Re\left(\frac{1}{1+iy}\right) -1\right) 
	= - \frac{\pi^2}{15\varepsilon} \frac{y^2}{1+y^2} \gg \varepsilon^{-2\delta}
\]
for $\varepsilon^{\frac12-\delta} \ll y \ll 1$. Thus, it follows that, for $\varepsilon^{\frac12-\delta} \ll y \ll 1$, we have, for all $L \in \mathbb{N}$,
\[
	\frac{\phi^{\frac12} }{5^{\frac14}}\log(\phi) e^{\frac{\pi^2}{15 z}-\frac{\pi^2}{15\varepsilon}} \ll \varepsilon^L.
\]
\end{proof}

\end{document}
